% from ega4-1-fr.tex
\oldpage[IV-1]{245}

De même, pour démontrer (vii), on peut se borner au cas où $X$ est affine~;
alors $X'$ est réunion finie d'ouverts affines $X'_i$ et $f(E')$ est réunion
des $f(E'\cap X'_i)$, si bien qu'on peut aussi supposer que $X'$ est affine,
en vertu de (ii)~; mais alors $E'=g(X'')$, où $g:X''\to X'$ est un morphisme
quasi-compact, en vertu de (ix), et l'on a $f(E')=f(g(X''))$, donc $f(E')$ est
pro-constructible puisque $g\circ f$ est quasi-compact.

On déduit de (vii) la démonstration de la seconde assertion de (v)~: en
effet, soit $E$ une partie localement fermée et rétrocompacte dans $X$, et
considérons un sous-préschéma de $X$ ayant $E$ pour espace sous-jacent (I,
5.2.1)~; l'injection canonique $j:E\to X$ est alors par hypothèse un
morphisme quasi-compact, et il résulte de (vii) que $E=j(E)$ est
pro-constructible dans $X$.

Enfin, pour prouver (viii), on peut encore supposer $X$ affine~; en outre, si
$(X'_\alpha)$ est un recouvrement ouvert affine de $X'$ tel que
$E'\cap X'_\alpha$ soit réunion de parties constructibles de $X'_\alpha$
((iii) et (iv)), $f(E')$ est réunion des $f(E'\cap X'_\alpha)$ et, en vertu
du théorème de Chevalley (1.8.4), chacune des $f(E'\cap X'_\alpha)$ est
réunion de parties constructibles de $X$, d'où la conclusion.
\end{proof}

\begin{env}[1.9.6]
\label{IV.1.9.6-fr}
\textbf{Exemples} (1.9.6). --- Toute partie \emph{finie} d'un préschéma $X$
est pro-constructible~: en effet, il suffit de considérer les parties
$\{x\}$ réduites à un seul point (1.9.5, (ii)) et $\{x\}$ est l'image de
$\operatorname{Spec}(k(x))$ par le morphisme canonique
$\operatorname{Spec}(k(x))\to X$, qui est quasi-compact~; d'où la conclusion
par (1.9.5, (vii)). Par contre, une partie $\{x\}$ réduite à un point n'est
pas nécessairement ind-constructible~; par exemple, soit $A$ un anneau
noethérien intègre ayant une infinité d'idéaux maximaux, et soit $x$ le point
générique de $X=\operatorname{Spec}(A)$~; si $\{x\}$ était ind-constructible
dans $X$, il serait constructible en vertu de (1.9.5, (iii)) et par suite
contiendrait un ouvert non vide de $X$ (0\textsubscript{III}, 9.2.2)~; mais
cela contredit l'hypothèse que $A$ possède une infinité d'idéaux maximaux, en
vertu du théorème d'Artin-Tate (0, 16.3.3).

Toute partie \emph{fermée} $Y$ d'un préschéma $X$ est pro-constructible, par
(1.9.5, (v)). Toute partie ouverte de $X$ est donc ind-constructible, mais
une partie ouverte $U$ de $X$ n'est pro-constructible que si elle est
\emph{rétrocompacte}, en vertu encore de (1.9.5, (v)).

Enfin, si $X$ est un préschéma \emph{noethérien}, donc quasi-séparé (1.2.8),
il résulte de (1.9.5, (iii)) et de (0\textsubscript{III}, 9.1.7) que les
parties ind-constructibles de $X$ sont les \emph{réunions de parties
localement fermées de $X$}.
\end{env}

\begin{theorem}[1.9.7]
\label{IV.1.9.7-fr}
Soient $X$ un préschéma quasi-compact, $E$ une partie ind-constructible de
$X$, $(F_\lambda)_{\lambda\in L}$ une famille de parties pro-constructibles
de $X$ telle que $\bigcap_{\lambda\in L}F_\lambda\subset E$~; alors il existe
une partie finie $J$ de $L$ telle que
$\bigcap_{\lambda\in J}F_\lambda\subset E$.
\end{theorem}

\begin{proof}
Notons que les ensembles $F=X-E$ et $F\cap F_\lambda$ sont
pro-constructibles, donc on est ramené au
\end{proof}

\begin{corollary}[1.9.8]
\label{IV.1.9.8-fr}
Soient $X$ un préschéma quasi-compact, $(F_\lambda)_{\lambda\in L}$ une
famille de parties pro-constructibles de $X$ dont l'intersection est vide~;
alors il existe une sous-famille finie $(F_\lambda)_{\lambda\in J}$ dont
l'intersection est vide.
\end{corollary}

\begin{proof}
Recouvrant $X$ par un nombre fini d'ouverts affines, et utilisant (1.9.5,
(iv)), on est ramené au cas où $X$ est affine~; en vertu de (1.9.5, (ix)), on
a $F_\lambda=f_\lambda(X'_\lambda)$,

